\documentclass[11pt]{article}
\usepackage[a4paper,margin=1in]{geometry}
\usepackage{fontspec}
\usepackage[boldfont=false]{xeCJK}
\setCJKmainfont{FandolSong-Regular.otf}
\usepackage{amsmath}
\usepackage{amssymb}
\usepackage{array}
\usepackage{booktabs}
\usepackage{graphicx}
\usepackage{adjustbox}
\usepackage{enumitem}
\setlist[itemize]{topsep=2pt,partopsep=0pt,parsep=0pt,itemsep=2pt,leftmargin=1.4em}
\setlist[enumerate]{topsep=2pt,partopsep=0pt,parsep=0pt,itemsep=2pt,leftmargin=1.6em}
\usepackage{parskip}
\usepackage{fvextra}
\fvset{breaklines=true,breakanywhere=true,fontsize=\small}
\setlength{\parindent}{0pt}

\begin{document}

\begin{center}
  {\LARGE\bfseries Atomic structure}\\[0.35em]
  {\large A Level Chemistry}
\end{center}
\vspace{0.6em}

\section*{What an atom is made of}

Everything is made of \textbf{atoms} 原子. An atom is mostly empty space. At its centre is a tiny, heavy \textbf{nucleus} 原子核. The nucleus holds two kinds of particle: \textbf{protons} 质子 and \textbf{neutrons} 中子. Around the nucleus, in the empty space, move the \textbf{electrons} 电子. The electrons stay in \textbf{shells} 壳层 — layers at set distances from the nucleus.

The nucleus is very small but holds almost all the mass. The electrons take up almost all the space but have almost no mass.

\subsection*{Relative charge and relative mass}

We compare the three particles using \textbf{relative charge} 相对电荷 and \textbf{relative mass} 相对质量. These are simple numbers, not real units.

\begin{center}
\adjustbox{max width=\linewidth}{%
\begin{tabular}{lll}
\toprule
\textbf{Particle} & \textbf{Relative charge} & \textbf{Relative mass} \\
\midrule
proton & $+1$ & $1$ \\
neutron & $0$ & $1$ \\
electron & $-1$ & $\tfrac{1}{1836}$ (about $0$) \\
\bottomrule
\end{tabular}}
\end{center}

A proton and a neutron have almost the same mass. An electron is about 1836 times lighter. The proton is positive, the electron is negative, and the neutron has no charge — it is \textbf{neutral} 中性.

\subsection*{Proton number and nucleon number}

Two numbers describe the nucleus:

\begin{itemize}
\item the \textbf{proton number} 质子数 (also called the \textbf{atomic number} 原子序数), symbol $Z$ — the number of protons.

\item the \textbf{nucleon number} 核子数 (also called the \textbf{mass number} 质量数), symbol $A$ — the total number of protons and neutrons. Protons and neutrons are both \textbf{nucleons} 核子.

\end{itemize}

So the number of neutrons is $A - Z$.

\subsection*{Counting particles in an atom or ion}

For a neutral atom, the number of electrons equals the number of protons, which equals $Z$.

An \textbf{ion} 离子 is an atom that has lost or gained electrons, so it has a charge:

\begin{itemize}
\item a positive ion has fewer electrons than protons.

\item a negative ion has more electrons than protons.

\end{itemize}

Example: $^{27}_{13}\text{Al}^{3+}$ has $13$ protons, $27 - 13 = 14$ neutrons, and $13 - 3 = 10$ electrons (it lost 3 electrons to become $3+$).

\subsection*{How mass and charge are spread out}

Almost all the mass sits in the nucleus, because protons and neutrons are heavy and electrons are very light. All the positive charge is in the nucleus (the protons). The negative charge is spread out in the shells (the electrons).

\subsection*{Beams of particles in an electric field}

Imagine beams of protons, neutrons and electrons moving at the same speed into an \textbf{electric field} 电场 between two charged plates:

\begin{itemize}
\item the proton beam bends towards the negative plate (protons are positive).

\item the electron beam bends the other way, towards the positive plate. It bends much more, because the electron is far lighter — the same force gives a bigger \textbf{deflection} 偏转 to a smaller mass.

\item the neutron beam goes straight through. It has no charge, so the field gives it no force.

\end{itemize}

\subsection*{Atomic radius and ionic radius}

The \textbf{atomic radius} 原子半径 is the size of an atom. The \textbf{ionic radius} 离子半径 is the size of an ion.

Across a \textbf{period} 周期 (left to right), the atomic radius gets \textbf{smaller}. The \textbf{nuclear charge} 核电荷 (the pull from the protons) rises, but the electrons go into the same outer shell, so the \textbf{shielding} 屏蔽 by inner shells stays about the same. The stronger pull draws the outer shell inwards.

Down a \textbf{group} 族 (top to bottom), the atomic radius gets \textbf{larger}. Each step down adds a new shell, so the outer electrons are further out and feel more shielding from the nucleus.

For ions:

\begin{itemize}
\item a positive ion (\textbf{cation} 阳离子) is smaller than its atom. It has lost its outer shell, and the electrons that remain feel a stronger pull each.

\item a negative ion (\textbf{anion} 阴离子) is larger than its atom. It has gained electrons, so there is more \textbf{repulsion} 排斥 between the electrons.

\item among ions that have the same number of electrons, the one with more protons is smaller.

\end{itemize}

\section*{Isotopes}

\textbf{Isotopes} 同位素 are atoms of the same element with the same number of protons but a different number of neutrons. So isotopes have the same proton number $Z$ but a different nucleon number $A$.

We write an isotope as $^{A}_{Z}\text{X}$: the nucleon number $A$ on top, the proton number $Z$ below. For example, chlorine has two main isotopes, $^{35}_{17}\text{Cl}$ and $^{37}_{17}\text{Cl}$.

\subsection*{Same chemical properties}

\textbf{Chemical properties} 化学性质 depend on the electrons, especially the outer electrons. Isotopes of one element have the same number of electrons arranged in the same way. So they react in exactly the same way — they have the same chemical properties.

\subsection*{Different physical properties}

Some \textbf{physical properties} 物理性质 depend on mass, so they differ between isotopes. A heavier isotope has more neutrons, so more mass, and therefore a higher \textbf{density} 密度. (The syllabus limits this difference to mass and density.)

\section*{Electrons, energy levels and orbitals}

Electrons are arranged in shells, sub-shells and orbitals.

\subsection*{Shells and the principal quantum number}

Each shell is labelled by the \textbf{principal quantum number} 主量子数 $n = 1, 2, 3, \dots$ A larger $n$ means a shell that is further from the nucleus and higher in energy.

\subsection*{Sub-shells and orbitals}

Each shell is split into \textbf{sub-shells} 亚层, named s, p and d. Each sub-shell is built from \textbf{orbitals} 轨道. An orbital is a small region that can hold up to two electrons.

\begin{center}
\adjustbox{max width=\linewidth}{%
\begin{tabular}{lll}
\toprule
\textbf{Sub-shell} & \textbf{Number of orbitals} & \textbf{Maximum electrons} \\
\midrule
s & 1 & 2 \\
p & 3 & 6 \\
d & 5 & 10 \\
\bottomrule
\end{tabular}}
\end{center}

So an s sub-shell holds 2 electrons, a p sub-shell holds 6, and a d sub-shell holds 10.

\subsection*{Order of increasing energy}

Electrons fill the lowest-energy sub-shell first. For the first three shells, plus 4s and 4p, the order of rising energy is:

\[
1\text{s} < 2\text{s} < 2\text{p} < 3\text{s} < 3\text{p} < 4\text{s} < 3\text{d} < 4\text{p}
\]

Notice the surprise: 4s is slightly lower in energy than 3d, so 4s fills first.

\subsection*{Electronic configuration}

The \textbf{electronic configuration} 电子排布 lists how many electrons are in each sub-shell. The lowest-energy arrangement is the \textbf{ground state} 基态.

For iron (Fe, $Z = 26$):

\[
1\text{s}^2\,2\text{s}^2\,2\text{p}^6\,3\text{s}^2\,3\text{p}^6\,3\text{d}^6\,4\text{s}^2
\]

You can write a shorthand using the nearest \textbf{noble gas} 稀有气体 in square brackets:

\[
[\text{Ar}]\,3\text{d}^6\,4\text{s}^2
\]

Here $[\text{Ar}]$ stands for the full configuration of argon.

For ions, you add or remove electrons. One key rule: when a transition metal forms a positive ion, it loses its 4s electrons \textbf{before} its 3d electrons. So $\text{Fe}^{3+}$ is $[\text{Ar}]\,3\text{d}^5$.

\subsection*{Electrons in boxes}

The \textbf{electrons in boxes} notation draws each orbital as a box and each electron as an arrow. Two electrons in the same orbital must point opposite ways, because each electron has a property called \textbf{spin} 自旋, and a shared orbital needs opposite spins.

Within a sub-shell, electrons fill empty orbitals one at a time, with parallel arrows, before any orbital gets a second electron. Spreading out like this keeps the electrons apart and lowers the repulsion between them.

\subsection*{Why the configuration takes this shape}

Electrons fill from low energy to high energy because that gives the most stable (lowest-energy) atom. Within a sub-shell they spread out singly first to reduce the repulsion between the negative electrons.

\subsection*{Shapes of s and p orbitals}

\begin{itemize}
\item an s orbital is a \textbf{sphere} 球形 centred on the nucleus.

\item a p orbital has two lobes, like a dumbbell, pointing along one axis. The three p orbitals point along three directions at right angles (the $x$, $y$ and $z$ axes).

\end{itemize}

\subsection*{Free radicals}

A \textbf{free radical} 自由基 is a species with one or more \textbf{unpaired electrons} 未成对电子. Free radicals are very reactive.

\section*{Ionisation energy}

\subsection*{First ionisation energy}

The \textbf{first ionisation energy} 第一电离能 (IE) is the energy needed to remove one electron from each atom in one mole of gaseous atoms, forming one mole of gaseous $+1$ ions.

We use gaseous atoms so there are no forces between the particles. The unit is $\text{kJ mol}^{-1}$. As an equation, for an element X:

\[
\text{X}(\text{g}) \rightarrow \text{X}^{+}(\text{g}) + \text{e}^{-}
\]

The $(\text{g})$ shows the species is a gas.

\subsection*{Successive ionisation energies}

After removing one electron, you can remove another. The second ionisation energy removes one electron from each $+1$ ion:

\[
\text{X}^{+}(\text{g}) \rightarrow \text{X}^{2+}(\text{g}) + \text{e}^{-}
\]

You can keep going. These are the \textbf{successive ionisation energies} 逐级电离能. Each is larger than the one before, because every electron is pulled away from a more positive ion.

\subsection*{What ionisation energy depends on}

Ionisation energy comes from the attraction between the positive nucleus and the outer electron. Three main factors set how strong that attraction is:

\begin{itemize}
\item \textbf{nuclear charge}: more protons pull the electrons more strongly, so the ionisation energy is higher.

\item \textbf{atomic radius}: the further the outer electron sits from the nucleus, the weaker the pull, so the ionisation energy is lower.

\item \textbf{shielding}: inner shells block some of the pull on the outer electron. More inner shells mean more shielding and a lower ionisation energy.

\end{itemize}

There is a smaller effect too — \textbf{spin-pair repulsion} 自旋成对排斥. When two electrons share one orbital, they push each other a little, so one is easier to remove.

\subsection*{Trends in first ionisation energy}

Across a period, the first ionisation energy generally \textbf{rises}. The nuclear charge grows while shielding stays about the same, so the outer electrons are held more tightly.

Down a group, the first ionisation energy \textbf{falls}. Lower elements have more shells, so more shielding and a larger radius, and the outer electron is easier to remove.

\subsection*{The dips are evidence for sub-shells}

The rise across a period is not smooth. Two small dips appear, and you should be able to explain both:

\begin{itemize}
\item \textbf{Group 2 to Group 13} (for example Mg to Al): the electron removed from Al comes from a 3p sub-shell, which is higher in energy than the full 3s sub-shell in Mg. A 3p electron is easier to remove, so the value dips.

\item \textbf{Group 15 to Group 16} (for example P to S): in S, one 3p orbital now holds a pair of electrons. Spin-pair repulsion makes one of them easier to remove, so the value dips.

\end{itemize}

These dips are evidence that sub-shells exist.

\subsection*{Successive ionisation energies are evidence for shells}

If you plot the successive ionisation energies of one element, the values rise, with \textbf{big jumps} at certain points. A big jump happens when the next electron must come from a shell closer to the nucleus.

Count how many electrons come off easily before the first big jump — that is the number of electrons in the outer shell, which tells you the \textbf{group} the element is in. You can also use the pattern to work out the electronic configuration and the position of the element in the Periodic Table.

\begin{center}
\adjustbox{max width=\linewidth}{%
\begin{tabular}{ll}
\toprule
\textbf{Electrons removed before the first big jump} & \textbf{Group} \\
\midrule
1 & Group 1 \\
2 & Group 2 \\
3 & Group 13 \\
\bottomrule
\end{tabular}}
\end{center}


\end{document}
